% BHU PYQ -- Auto-generated & error-corrected
\documentclass[11pt,a4paper]{article}

\usepackage[a4paper,top=2.5cm,bottom=2.5cm,left=2.8cm,right=2.8cm,headheight=14pt]{geometry}
\usepackage{amsmath,amssymb}
\usepackage[T1]{fontenc}
\usepackage[utf8]{inputenc}
\usepackage{lmodern}
\usepackage{microtype}
\usepackage{fancyhdr}
\usepackage{enumitem}
\usepackage{hyperref}
\usepackage{lastpage}
\usepackage{parskip}

\hypersetup{colorlinks=true,urlcolor=black,linkcolor=black,
  pdftitle={STA-201: Descriptive Statistics and Distribution Theory -- BHU PYQ},pdfauthor={Banaras Hindu University}}

\pagestyle{fancy}\fancyhf{}
\renewcommand{\headrulewidth}{0.4pt}\renewcommand{\footrulewidth}{0.4pt}
\fancyhead[L]{\small   \textit{Banaras Hindu University}}
\fancyhead[R]{\small   \textit{STA-201: Descriptive Statistics and Distribu}}
\fancyfoot[C]{\small Page  \thepage\ of \pageref{LastPage}}


\newlist{parts}{enumerate}{2}
\setlist[parts,1]{label=(\alph*),leftmargin=*,itemsep=5pt,topsep=3pt}
\setlist[parts,2]{label=(\roman*),leftmargin=*,itemsep=3pt,topsep=2pt}

\newcommand{\pts}[1]{\hfill   \textit{[\#1]}}


\begin{document}


\begin{center}
  {\large\bfseries BANARAS HINDU UNIVERSITY}\\[2pt]
  {\normalsize B.A. (Hons.) Semester II Examination, 2022-23}\\[6pt]
  \rule{0.8\linewidth}{0.5pt}\\[6pt]
  {\large\bfseries Paper No. STA-201: Descriptive Statistics and Distribution Theory}\\[4pt]
  {\normalsize Time: 3 Hours\hfill Maximum Marks: 70}
\end{center}
\rule{\linewidth}{0.4pt}
\smallskip



\noindent   \textit{Instructions: Answer   \textbf{five} questions including Question~1 which is compulsory. Symbols have their usual meanings.}
\bigskip

\medskip


\noindent   \textbf{Question 1.}

\begin{parts}
  \item Define mathematical expectation. If $X$ is a random variable, then prove that $|E(X)| \le E(|X|)$. \pts{7}
  \item Write a note on linear and non-linear correlation. Also show that Karl Pearson's coefficient of correlation lies between -1 and +1. \pts{7}
\end{parts}
\medskip\rule{\linewidth}{0.2pt}

\medskip


\noindent   \textbf{Question 2.}

\begin{parts}
  \item If $X_1$, $X_2$, and $X_3$ are uncorrelated variables each having the same variance, obtain the correlation coefficient between $X_1 + X_2$ and $X_2 + X_3$. \pts{7}
  \item A random variable $X$ can assume any positive integral value $n$ with a probability proportional to $1/3^n$. Find the expectation of $X$. \pts{7}
\end{parts}
\medskip\rule{\linewidth}{0.2pt}

\medskip


\noindent   \textbf{Question 3.}

\begin{parts}
  \item Define intra-class correlation and derive the expression for the intra-class correlation coefficient. \pts{7}
  \item Define convergence in probability. Let $X_1, X_2, \dots, X_n$ be i.i.d. random variables with probability density function $f(x) = e^{-x}, x > 0$. Show that $Y_n  o 0$ in probability, where $Y_n = \min(X_1, X_2, \dots, X_n)$. \pts{7}
\end{parts}
\medskip\rule{\linewidth}{0.2pt}

\medskip


\noindent   \textbf{Question 4.}

\begin{parts}
  \item For the joint distribution of two-dimensional random variables $(X, Y)$ given by:
  \[
  f(x, y) = e^{-(x+y)}, \quad x > 0, \ y > 0
  \]
  Show that the characteristic function of $X+Y$ is equal to the product of the characteristic functions of $X$ and $Y$. \pts{7}
  \item Find the value of the correlation coefficient for the data set $(X, Y)$ in which $Y$ has the same constant value corresponding to any value of $X$. \pts{7}
\end{parts}
\medskip\rule{\linewidth}{0.2pt}

\medskip


\noindent   \textbf{Question 5.}

\begin{parts}
  \item For two variables $X$ and $Y$ with the same mean, the two regression equations are $y = ax + b$ and $x = \alpha y + \beta$. Find the mean and show that $\frac{b}{\beta} = \frac{1-a}{1-\alpha}$. \pts{7}
  \item Define the consistency of data. If data has three attributes, give the conditions of the consistency of the data. \pts{7}
\end{parts}
\medskip\rule{\linewidth}{0.2pt}

\medskip


\noindent   \textbf{Question 6.}

\begin{parts}
  \item Derive the Poisson distribution as a limiting case of the Binomial distribution. Also obtain the mode of the Poisson distribution. \pts{10}
  \item Write a note on the lack of memory property of the geometric distribution. \pts{4}
\end{parts}
\medskip\rule{\linewidth}{0.2pt}

\medskip


\noindent   \textbf{Question 7.}

\begin{parts}
  \item Define the gamma distribution with one as well as two parameters. Derive the moment generating function for the two-parameter gamma distribution. \pts{10}
  \item If $X$ is uniformly distributed with mean $1$ and variance $4/3$, find $P(X < 0)$. \pts{4}
\end{parts}
\medskip\rule{\linewidth}{0.2pt}

\medskip


\noindent   \textbf{Question 8.}
Write a note on any   \textbf{three} of the following: \pts{14}

\begin{parts}
  \item Rank Correlation
  \item Normal Distribution
  \item Association of Attributes
  \item Chebyshev's Inequality
\end{parts}

\vfill
\rule{\linewidth}{0.4pt}

\begin{center}\small
  \href{https://bhu-exam.vercel.app/}{Click here to visit our website}\\[4pt]\href{https://me-aryan.vercel.app/}{Click here to visit our contributor}\\[4pt]\href{https://quashan-me.vercel.app/}{Click here to visit our contributor}
\end{center}

\end{document}